% ---
% id: protocol-compiler/def:compile
% kind: def
% title: "The compilation map"
% summary: >-
%   The compilation map fills the declaration from the specification, passing
%   declared data through and constructing only the raw-key map, the index map
%   and the weighted families. In the general format it also assembles the
%   factor list.
% prominence: standard
% status: implemented
% uses: [def:declaration, def:homog, def:target, def:exact, transmitter/def:source-spec]
% mentions: [def:weighted, thm:sifting-and-key, def:gate-outcome-table, def:eligible, def:validator, def:spec, def:chaincond, reductions/def:joint-family, transmitter/rem:da1, transmitter/req:marginal-fill, foundations/req:declared-order, req:pass-through, req:derived-objects, req:validator-refusals]
% assumes: []
% refs: []
% provenance: newly-derived
% implements: certiqs_key.protocol_compiler.compile:compile_specification
% lean:
% notes: "The edge to def:eligible is a mention: the gate is decided on the emitted data after this map has run, so the map does not depend on it."
% ---

The compilation map assembles the declaration from the specification. It
passes the \tref{declared data} through, weights a measurement family where a public
choice is presented, and constructs only the raw-key map, the index map and
the weighted families. The list below fills the \tref{basic format}; the
paragraphs on the general format extend it slot for slot. That the engine
passes the declared data through unchanged, constructs no right-hand side, and
forms no product of two factors' families, no \tref{lifted operator} and no
other derived object of the general format as data is
Requirements~\ref{req:pass-through} and~\ref{req:derived-objects}; the
refusals the validator issues on the factor list and on Alice's slots are
Requirement~\ref{req:validator-refusals}'s.

At $d_A=1$ a trusted-source request may leave the marginal slot empty; the
transmitter then fills it with $[1]$ and records the fill
(\xref{transmitter/req:marginal-fill}); why the two
source modes impose the same condition there is \xref{transmitter/rem:da1}. A
certificate that is present and fails any of the five conditions of gate (G2)
is refused at every $d_A$, row~3 of the \tref{gate outcome table} of
Definition~\ref{def:gate-outcome-table}. The order in which the binding rule below has
condition~(S) decided on a supplied family is the order that family was
declared in and also the order the consuming engine re-accumulates in (\xref{foundations/req:declared-order}): one order
at all three places.

\begin{definition}\label{def:compile}
The \term[The map from a protocol specification to the protocol declaration it fills, slot by slot.]{compilation map} is the map $\Spec\mapsto\Decl$ that fills
\eqref{eq:declaration} as follows.
\begin{itemize}[leftmargin=1.6em]
\item \emph{Alice's system and measurement.} $A$, $d_A$ and $\{M_x\}$ from
  $\Spec_{\mathrm{src}}$, weighted by Definition~\ref{def:weighted} where a
  choice family is presented.
\item \emph{Alice's marginal.} $\tau_A$ is what the transmitter artefact's
  marginal slot holds: in trusted-source mode the matrix in the slot --- at
  $d_A=1$ the value $[1]$ the transmitter filled the slot with, the unique
  element of $\Seq{A}$ --- and
  $\tau_A=\nomarg$ where the slot is empty. The mode is read off that slot and
  off nothing else. A certificate the artefact carries is consumed by gate
  (G2) of Definition~\ref{def:eligible}.
\item \emph{The \tref{measured system} and its family.} $(\Qsys,\SecSet,(d_s)_s)$
  and $\{F_v\}$ from the last link of $\Chrec$. Where $\Spec_{\mathrm{rec}}$
  declares a setting-dependent public choice, the family is weighted
  after the chain by the setting-dependent form of
  Definition~\ref{def:weighted}: the emitted family is then
  $N_{v,f}=P(f\,|\,t(v))\,F_v$ over the support of $P$, the emitted \tref[outcome alphabet]{outcome
  alphabet} is that support, and the setting map $t$ through which $P$ is read
  is the class map of the \tref{joint alphabet} of the chain's \tref{producer} where
  the chain was emitted by the reductions module
  (\xref{reductions/def:joint-family}\,(a)), and is declared with the family
  otherwise.
\item \emph{The public maps.} $\Party$, $a_A$, $a_B$, $\phi$, $\hold$
  verbatim from $\Spec_{\mathrm{pub}}$, declared on the emitted alphabets.
\item \emph{The constraint list.} $\CList$ is the concatenation of
  $\Spec_{\mathrm{con}}$'s entries with the entries transported by the links
  of $\Chsrc$ and then those of $\Chrec$, each segment in its own order, each
  entry with the datum its right-hand-side slot holds and with the tag that
  datum's type determines (Definition~\ref{def:homog}).
\item \emph{The raw-key map.}
\begin{equation}\label{eq:rdef}
  r(x,v)\coloneqq
  \begin{cases}
    r^{0}_{\hold}\big(\rd{\hold}(x,v),\,a(x,v)\big),
      & \varsigma(a(x,v))=\mathrm{pass},\\[0.3ex]
    \perp, & \varsigma(a(x,v))=\mathrm{fail},
  \end{cases}
  \qquad\text{with }d_Z\coloneqq|\Zval| .
\end{equation}
\end{itemize}

\emph{The index map.} Alongside $\Decl$ the compiler emits the
\term[The emitted choice of a bijection from each alphabet and from the sector set onto an initial segment of the non-negative integers, through which every emitted list is indexed.]{index map}
$\idxm$, the label-to-index map: a choice, for each of $\Xset$, $\Vset$,
$\SecSet$, $\Oset{\mathsf A}$, $\Oset{\mathsf B}$, $\CAset$, $\CBset$,
$\Chatset$ and $\Zset$, of a bijection onto an initial segment of
$\mathbb{Z}_{\ge0}$; every emitted list, matrix and function table is indexed
through $\idxm$, and every semantic label is carried as metadata only. Its
component on $\SecSet$ is the total order the sector map $\Smap$ already
fixes, and not a second choice: one order on the sectors, required by the last
clause of D(i) (Definition~\ref{def:target}) and decided by V1. And $\idxm$ is
what fixes the declared summation order of Definition~\ref{def:exact} on every
alphabet it indexes; it indexes no basis of a matrix, a trace being
accumulated in the row order of the emitted array.

\emph{The binding rule.} On the remaining alphabets the choice is the
compiler's, under one binding rule. Where a supplied family arrives with a
declared order on the alphabet that indexes it --- as the target POVM of the
last link of $\Chrec$ does, its \tref{declarer} having established its own
completeness sum in that order --- the compiler takes \emph{that} order as
$\idxm$'s component there. So the order in which condition~(S) of
Definition~\ref{def:exact} is decided on a supplied family here is the order
that family was declared in, and no family is refused here for an ordering
its declarer never chose. Where no order arrives with an alphabet the choice is free, and (S) is
decided in the order the compiler chose and emits. This rule prescribes
nothing to the declarer of a family: it constrains only which bijection this
module picks.

Where a choice-weighted presentation is resolved, this same rule settles the
order the decision on each conditional family is read in, and adds nothing to
it. The family the weighting produces is indexed by one alphabet --- $\Xset$
where the presentation is Alice's, the physical receiver alphabet where it is
a measuring party's, which in \tref[Band A]{Band~A} is the emitted $\Vset$ --- and where a declared order
arrives with that alphabet the rule above takes \emph{it}, here as everywhere:
$\Xset$ arrives ordered with the source part
(Definition~\ref{def:spec}\,(a)), and the physical alphabet with the order the
fibre sum of Definition~\ref{def:chaincond}\,(C4) is accumulated in.
Condition~(S) on the conditional family of a choice value is then read in the
\emph{restriction of that order} to the elements that choice value
contributes, which is a total order on $\Uset_t$ however the record listed
it; the two conditions are independent, only the order restricting. Where no order arrives with that alphabet the
free choice is made as Definition~\ref{def:weighted} says, by the record's own
listing, and it is emitted. No declarer's order is overridden and no
presentation is refused for an order. Where the setting-dependent weighting is applied,
the emitted alphabet is the support of $P$ and its order is the product order
of Definition~\ref{def:weighted}: $v$ in its declarer's declared order and,
within each $v$, $f$ in the declared order of the \tref{choice alphabet}, the
pairs off the support omitted; that is the order (S) is decided in on the
emitted family and the order $\idxm$ emits on it.

\emph{The general format.} For a specification presented as
Definition~\ref{def:spec}'s (a$'$)--(e$'$) the map fills the general format
of Definition~\ref{def:declaration}\,(DX1)--(DX6) as follows. In the
\tref{basic format} --- one Alice sector, one measured factor, POVM form, two
messages neither reading the other and no party-local entry --- every clause
below is the clause of the same subject above, slot for slot, and the map
emits what those clauses emit.

\emph{The factor list} (DX1) is assembled and not declared as a whole: the
factor tagged $A$ is the Alice sector map the transmitter artefact declares,
where it declares one, and it is the first of the list; where the artefact
declares none, $\mathbb{T}=\{t_0\}$ and no $A$ factor is emitted. The factors
tagged $\Qsys$ are the entries of $\Spec_{\mathrm{rec}}$ in their declared
order, each carrying that entry's party label and the sector map of the last
target descriptor of that entry's receiver segment, which
Definition~\ref{def:chaincond}\,(C1) decides per factor. The one block kind
is the kind those maps carry, and the factors' agreement in it is decided by
the same four-component comparison V1 makes, taken on the factor list
(Definition~\ref{def:validator}, V1).
The product sector map, its order, the product alphabet and the
\tref[joint sector]{joint sectors} are derived from that list by (DX1) and (DX3) and are
declared, formed and emitted by no one: what is emitted is the factor list and
the rules.

\emph{Alice's slots.} $\tau_A$ is what the artefact's marginal slot holds,
the block list along $\mathbb{T}$ included; Alice's measurement slot is the
artefact's, in the form the artefact declares --- the POVM slot or the
ensemble slot of (DX4) --- weighted by Definition~\ref{def:weighted} in that
form where a choice family is presented. The mode is read off the
marginal slot as above.

\emph{Per measured factor.} $(\Qsys_f,\SecSet^{f},(d_{f,s_f})_{s_f})$ and
$\{F^{f}_{v_f}\}_{v_f\in\Vset^{f}}$ from the last link of that factor's
receiver segment and, where that factor's entry of $\Spec_{\mathrm{rec}}$
declares a setting-dependent public choice, weighted after that segment by the
setting-dependent form of Definition~\ref{def:weighted}, so that factor $f$'s
emitted family is $P_f(\,\cdot\mid\cdot\,)$ times that factor's delivered
family over the support of $P_f$, and factor $f$'s emitted alphabet
$\Vset^{f}$ is that support in its product order. Each factor's weighting is
applied to its own family and to nothing else. No product of two factors'
families is formed by this map, whether emitted or handed on; so the subject of the
representability decision of Definition~\ref{def:weighted} is, as in the basic
format, the products \emph{that definition's} two constructions form, and no
product block \eqref{eq:productfamily} is a product this map forms.

\emph{The public part and the raw key.} $\Party$, $\phi$, $\hold$ verbatim
from $\Spec_{\mathrm{pub}}$ and the \tref[transcript protocol]{transcript} verbatim: the sender,
the alphabet and the table of each message, declared on the emitted
alphabets. The joint announcement function $a$ is not a slot: it is the
composition \eqref{eq:transcript} of those tables, and $r$ is \eqref{eq:rdef}
with \emph{that} $a$, which is what Theorem~\ref{thm:sifting-and-key}\,(iv), (vi) is
stated of. $\CList$ is the concatenation above, each entry carrying in
addition its \tref{scope}, and each with the datum its right-hand-side slot
holds and the tag that datum's type determines; for a party-local entry that
is the tag field it arrives with beside a real scalar
(Definition~\ref{def:homog}).

\emph{The index map in the general format.} $\idxm$ has one component per
factor alphabet $\Vset^{f}$ and one per factor \tref{sector set} $\SecSet^{f}$, and,
where an $A$ factor is emitted, one on $\mathbb{T}$, which is the total order
that factor's own map fixes. Its components on $\Vset$ and on $\SecSet$ are
the product orders (DX1) and (DX3) derive from the factors' orders and are not
further choices, as its component on $\SecSet$ above is not; the joint sectors
are ordered with $t$ slowest by (DX1). The binding rule above is applied per
factor: where a factor's target family arrives with a declared order on its
alphabet, the compiler takes that order there. So condition~(S) of
Definition~\ref{def:exact} is decided per factor, in that factor's arriving
order, and per Alice sector on Alice's data; no sum over a product alphabet is
accumulated by this map or by its consumer.
\end{definition}
